\section[A basis for~$\Delta $ that involves~$\Omega$]{A basis for $\boldsymbol{\Delta}$ that involves~$\boldsymbol{\Omega}$}

 In Theorem~\ref{thm:basis} we displayed a basis for
$\Delta $. In this section we display a related
basis for $\Delta $ that involves the Casimir element~$\Omega$.
In the next section we will use the related basis
to describe  the center~$Z(\Delta )$.

 Recall the f\/iltration $\lbrace \Delta_n \rbrace_{n=0}^\infty$
of $\Delta $ from Def\/inition~\ref{def:ui}.


\begin{Lemma}
\label{lem:omegapower}
For all integers $\ell \geq 1$
the following hold:
\begin{enumerate}\itemsep=0pt
\item[$(i)$] $\Omega^\ell \in \Delta_{3 \ell}$.
\item[$(ii)$] $\Omega^\ell -q^{\ell^2}A^\ell B^\ell C^\ell  \in
\Delta_{3\ell-1}$.
\end{enumerate}
\end{Lemma}

\begin{proof}
Consider the expression for $\Omega$ from the
f\/irst displayed line of
Lemma~\ref{lem:six}.
The term $qABC$ is in $\Delta_3$
and the remaining terms are in $\Delta_2$. Therefore
$\Omega \in \Delta_3$ and $\Omega-qABC \in \Delta_2$. By this and
Lemma~\ref{lem:filt}$(iii)$ we f\/ind
$\Omega^\ell \in \Delta_{3\ell}$ and
$\Omega^\ell - q^\ell (ABC)^\ell
\in \Delta_{3\ell-1}$.
Using Lemma~\ref{lem:almostcom}  we obtain
$(ABC)^\ell -q^{\ell (\ell-1)}A^\ell B^\ell C^\ell \in \Delta_{3\ell -1}$.
By these comments
$\Omega^\ell -q^{\ell^2}A^\ell B^\ell C^\ell  \in \Delta_{3\ell-1}$.
\end{proof}

\begin{Lemma}
\label{lem:basis2c}
For all integers $n\geq 1$
the following is a basis
for a complement of
$\Delta_{n-1}$ in $\Delta_n$:
\begin{gather*}
 A^iB^jC^k\Omega^\ell \alpha^r \beta^s \gamma^t, \qquad
i,j,k,\ell, r,s,t\geq 0,
\qquad ijk=0,
 \qquad
i+j+k+3\ell+r+s+t = n.
\end{gather*}
\end{Lemma}

\begin{proof}
Let $\mathbb I_n$ denote the set consisting of the 6-tuples
$(i,j,k,r,s,t)$ of nonnegative integers whose sum is~$n$.
By Corollary~\ref{lem:com} the
following
is a basis for a complement of
$\Delta_{n-1}$ in $\Delta_n$:
\begin{gather*}
A^iB^jC^k\alpha^r\beta^s\gamma^t, \qquad
(i,j,k,r,s,t) \in {\mathbb I}_n.
\end{gather*}
Let $\mathbb J_n$ denote the set consisting of the 7-tuples
$(i,j,k,\ell,r,s,t)$ of nonnegative integers such that
$ijk=0$ and
$i+j+k+3\ell+r+s+t = n$.
Observe that the map
\begin{gather*}
\mathbb J_n  \to  \mathbb I_n,
\qquad
(i,j,k,\ell,r,s,t)  \mapsto   (i+\ell,j+\ell,k+\ell,r,s,t)
\end{gather*}
is a bijection.
Suppose we are given $(i,j,k,\ell,r,s,t) \in \mathbb J_n$.
By
Lemma \ref{lem:omegapower},
$\Delta_{n-1}$ contains
\begin{gather*}
A^iB^jC^k\Omega^\ell \alpha^r \beta^s \gamma^t -
q^{\ell^2} A^iB^jC^k A^\ell B^\ell C^\ell \alpha^r \beta^s \gamma^t.
\end{gather*}
By Lemma~\ref{lem:almostcom},  $\Delta_{n-1}$ contains
\begin{gather*}
A^iB^jC^k A^\ell B^\ell C^\ell \alpha^r \beta^s \gamma^t
-
q^{2j \ell} A^{i+\ell}B^{j+\ell}C^{k+\ell}  \alpha^r \beta^s \gamma^t.
\end{gather*}
Therefore $\Delta_{n-1}$ contains
\begin{gather*}
A^iB^jC^k\Omega^\ell \alpha^r \beta^s \gamma^t -
q^{\ell(2j+ \ell)} A^{i+\ell}B^{j+\ell}C^{k+\ell}  \alpha^r \beta^s \gamma^t.
\end{gather*}
By these comments the following is a basis for a complement of
$\Delta_{n-1}$ in $\Delta_n$:
\begin{gather*}
A^iB^jC^k\Omega^\ell \alpha^r\beta^s\gamma^t, \qquad
(i,j,k,\ell,r,s,t) \in \mathbb J_n.
\end{gather*}
The result follows.
\end{proof}

\begin{Note}
Pick an integer $n \geq 1$.
In
Corollary
\ref{lem:com} and
Lemma
\ref{lem:basis2c}
we mentioned a
complement of $\Delta_{n-1}$ in $\Delta_n$. These
complements are not the same
in general.
\end{Note}


\begin{Proposition}
\label{lem:basis2n}
For all integers $n\geq 0$
the following is a basis
for
the $\mathbb F$-vector space
$\Delta_n$:
\begin{gather*}
 A^iB^jC^k\Omega^\ell \alpha^r \beta^s \gamma^t, \qquad
i,j,k,\ell, r,s,t\geq 0,
\qquad ijk=0,
 \qquad
i+j+k+3\ell+r+s+t \leq  n.
\end{gather*}
\end{Proposition}

\begin{proof}
By Lemma~\ref{lem:basis2c} and
$\Delta_0=\mathbb F 1$.
\end{proof}




\begin{Theorem}
\label{lem:basis2}
The following is a basis for the $\mathbb F$-vector space
$\Delta$:
\begin{gather*}
A^iB^jC^k\Omega^\ell \alpha^r \beta^s \gamma^t, \qquad
i,j,k,\ell, r,s,t\geq 0,
 \qquad ijk=0.
\end{gather*}
\end{Theorem}

\begin{proof}
Combine
Lemma~\ref{lem:filt}$(ii)$
and
Proposition~\ref{lem:basis2n}.
\end{proof}







\section[The center $Z(\Delta)$]{The center $\boldsymbol{Z(\Delta)}$}

In this section we give a detailed description of
the center $Z(\Delta)$, under the assumption that $q$ is not a root of
 unity.
For such $q$ we show that $Z(\Delta)$ is generated by
$\Omega$, $\alpha$, $\beta$, $\gamma$ and
 isomorphic to a polynomial algebra in four
variables.

  Recall the commutator
 $\lbrack r,s\rbrack=rs-sr$.


\begin{Lemma}
\label{lem:ijk}
Let
$i$, $j$, $k$
denote nonnegative integers.
Then $\Delta_{i+j+k}$ contains each of the fol\-lo\-wing:
\begin{gather}
\label{eq:1}
 \big\lbrack A,A^iB^jC^k\big\rbrack - \big(1-q^{2j-2k}\big)A^{i+1}B^jC^k,
\\
\label{eq:2}
 \big\lbrack B,A^iB^jC^k\big\rbrack - \big(q^{2i}-q^{2k}\big)A^iB^{j+1}C^k,
\\
\label{eq:3}
 \big\lbrack C,A^iB^jC^k\big\rbrack - \big(q^{2j-2i}-1\big)A^iB^jC^{k+1}.
\end{gather}
\end{Lemma}

\begin{proof}
Concerning (\ref{eq:1}), observe
\begin{gather*}
\big\lbrack A,A^iB^jC^k\big\rbrack = A^{i+1}B^jC^k -A^iB^jC^kA.
\end{gather*}
By Lemma
\ref{lem:almostcom} $\Delta_{i+j+k}$ contains
\begin{gather*}
A^iB^jC^kA - q^{2j-2k}A^{i+1}B^jC^k.
\end{gather*}
By these comments
$\Delta_{i+j+k}$ contains
(\ref{eq:1}).
One similarly f\/inds that
$\Delta_{i+j+k}$ contains
(\ref{eq:2}),
(\ref{eq:3}).
\end{proof}


\begin{Theorem}
\label{thm:centerbasis}
Assume that $q$ is not a root of unity. Then the following is a basis for
the $\mathbb F$-vector space~$Z(\Delta)$.
\begin{gather}
\Omega^\ell \alpha^r \beta^s\gamma^t, \qquad  \ell,r,s,t \geq 0.
\label{eq:cbasis}
\end{gather}
\end{Theorem}

\begin{proof}
Abbreviate $Z=Z(\Delta)$.
The elements~(\ref{eq:cbasis})
are linearly independent by
Theorem~\ref{lem:basis2},
so it suf\/f\/ices to
show that they span~$Z$.
Let~$Z'$ denote the subspace of~$\Delta$ spanned by~(\ref{eq:cbasis}), and note that
$Z' \subseteq Z$. To show
$Z'= Z$, we assume that $Z'$ is properly contained in
$Z$ and get a contradiction.
Def\/ine the set $E=Z\backslash Z'$
and note $E\not=\varnothing$.
We have
$\Delta_0=\mathbb F 1
\subseteq Z'$ so
$E \cap \Delta_0 =\varnothing$.
By this and
Lemma~\ref{lem:filt}$(i)$,$(ii)$ there exists a unique integer
$n\geq 1$ such that
$E\cap \Delta_n \not=\varnothing$ and
$E\cap \Delta_{n-1} =\varnothing$.
Fix $u \in E\cap \Delta_n$.
Let
$S=S(u)$ denote
the set of 6-tuples $(i,j,k,r,s,t)$ of nonnegative integers
whose sum is $n$ and
 $A^iB^jC^k\alpha^r\beta^s\gamma^t$
contributes to $u$ in
the sense of
Def\/inition~\ref{def:contrib}.
By~(\ref{eq:bilsum}) and Corollary~\ref{lem:com},
$\Delta_{n-1}$ contains
\begin{gather}
\label{eq:dif}
u- \sum_{(i,j,k,r,s,t) \in S}
\big( u,
A^iB^jC^k\alpha^r \beta^s \gamma^t \big)
A^iB^jC^k\alpha^r \beta^s \gamma^t.
\end{gather}
We are going to show that
$i=j=k$ for all
$(i,j,k,r,s,t) \in S$.
To this end
we f\/irst claim that $j=k$ for all
$(i,j,k,r,s,t) \in S$.
To prove the claim,
take the commutator of $A$ with~(\ref{eq:dif}) and evaluate the
result using the following facts.
By construction $u \in E \subseteq Z$
so
  $\lbrack A,u\rbrack=0$.
By Lemma~\ref{lem:filt}$(iii)$
$A\Delta_{n-1} \subseteq  \Delta_n $
and
$\Delta_{n-1}A \subseteq \Delta_n$
so
 $\lbrack A,\Delta_{n-1}\rbrack \subseteq \Delta_n$.
Moreover
 each of~$\alpha$,~$\beta$,~$\gamma$
 is central.
The above evaluation shows that~$\Delta_n$ contains
\begin{gather*}
\sum_{(i,j,k,r,s,t) \in S}
\big( u, A^iB^jC^k
\alpha^r \beta^s \gamma^t \big)
\big\lbrack A, A^iB^jC^k \big\rbrack
\alpha^r \beta^s \gamma^t.
\end{gather*}
Pick
$(i,j,k,r,s,t) \in S$.
By
Lemma~\ref{lem:filt}$(iii)$ and
Lemma~\ref{lem:ijk}, $\Delta_n$ contains
\begin{gather*}
\big\lbrack A, A^iB^jC^k \big\rbrack
\alpha^r \beta^s \gamma^t
-
\big(1-q^{2j-2k}\big)  A^{i+1}B^jC^k
\alpha^r \beta^s \gamma^t.
\end{gather*}
By the above comments
$\Delta_n$ contains
\begin{gather*}
\sum_{
(i,j,k,r,s,t) \in S}
\big(u, A^iB^jC^k
\alpha^r \beta^s \gamma^t \big)
\big(1-q^{2j-2k}\big)  A^{i+1}B^jC^k
\alpha^r \beta^s \gamma^t.
\end{gather*}
For all
$(i,j,k,r,s,t) \in S$
the element
$A^{i+1}B^jC^k
\alpha^r \beta^s \gamma^t$
is contained
in the basis for the complement of $\Delta_n$ in $\Delta_{n+1}$
given in Corollary~\ref{lem:com}. Therefore
\begin{gather*}
\big(u, A^iB^jC^k
\alpha^r \beta^s \gamma^t\big)
\big(1-q^{2j-2k}\big)  = 0
\qquad
\forall \;(i,j,k,r,s,t) \in S.
\end{gather*}
By the def\/inition of $S$ we have
$( u, A^iB^jC^k
\alpha^r \beta^s \gamma^t )\not=0$
 for
all
$(i,j,k,r,s,t) \in S$.
Therefore
$1-q^{2j-2k} = 0 $ for all
$(i,j,k,r,s,t) \in S$.
The scalar $q$ is not a root of unity
so $j=k$ for all
$(i,j,k,r,s,t) \in S$.
The claim is proved.
We next claim that
$i=j$ for all
$(i,j,k,r,s,t) \in S$.
This claim is proved like the previous one,
 except that
as we begin the argument below~(\ref{eq:dif}),
we use~$C$ instead of~$A$ in the commutator.
By the two claims
$i=j=k$ for all
$(i,j,k,r,s,t) \in S$.
In light of this we revisit the assertion above~(\ref{eq:dif}) and conclude that
$\Delta_{n-1}$ contains
\begin{gather*}
u-\sum_{
(i,i,i,r,s,t) \in S}
\big(u, A^iB^iC^i
\alpha^r \beta^s \gamma^t \big)
A^{i}B^iC^i
\alpha^r \beta^s \gamma^t.
\end{gather*}
Pick
$(i,i,i,r,s,t) \in S$.
By
Lemma~\ref{lem:filt}$(iii)$ and
Lemma~\ref{lem:omegapower},
$\Delta_{n-1}$ contains
\begin{gather*}
A^iB^iC^i
\alpha^r \beta^s \gamma^t
-q^{-i^2}\Omega^i \alpha^r \beta^s \gamma^t.
\end{gather*}
By these comments
$\Delta_{n-1}$ contains
\begin{gather*}
\label{eq:summ}
u-
\sum_{
(i,i,i,r,s,t) \in S}
q^{-i^2}
\big(u, A^iB^iC^i
\alpha^r \beta^s \gamma^t  \big)
\Omega^i \alpha^r \beta^s \gamma^t.
\end{gather*}
In the above expression let
$\psi$ denote the main sum,
so that
$u-\psi \in \Delta_{n-1}$.
Observe
$\psi \in Z' \subseteq Z$.
Recall $u \in E=Z\backslash Z'$ so
$u \in Z$ and $u \not\in Z'$. By these comments
$u - \psi \in Z$ and
$u - \psi \not\in Z'$.
Therefore $u - \psi \in E$ so
 $u - \psi \in E\cap \Delta_{n-1}$.
This contradicts
 $E\cap \Delta_{n-1}=\varnothing $ so
$Z=Z'$. The result follows.
\end{proof}
